% =============================================================================
% Section 5: Phase 1 — LoRA Encoder Adaptation
% =============================================================================
\section{Phase 1: LoRA Encoder Adaptation}
\label{sec:lora}

% ---------------------------------------------------------------------------
% 5.1 Adaptation Strategy
% ---------------------------------------------------------------------------
\subsection{Adaptation Strategy}
\label{sec:lora:strategy}

We adapt the frozen BrainSegFounder encoder to the meningioma domain via LoRA
(\cref{sec:bg:lora}), injecting low-rank adapter matrices into the
self-attention projections of Stages~3--4. Stages~0--2 remain frozen to
preserve low-level anatomical features learned from the UK~Biobank SSL
pretraining. Segmentation on BraTS-MEN 2024 serves as the proxy task.

\subsubsection{Target Modules}

LoRA adapters are injected into the $\mathbf{Q}$, $\mathbf{K}$, $\mathbf{V}$
projection matrices within the multi-head self-attention blocks of Stages~3
and~4:
\begin{align}
  &\texttt{swinViT.layers3.0.blocks.\{0,1\}.attn.qkv.weight}, \\
  &\texttt{swinViT.layers4.0.blocks.\{0,1\}.attn.qkv.weight}.
\end{align}
This yields 4 adapted layers, each with rank-$r$ decompositions for Q, K,
and V (combined as a single QKV projection in the Swin implementation). The
total number of trainable LoRA parameters is approximately $197\text{K}$ for
$r=8$, compared to 62M total model parameters ($<0.32\%$).

\subsubsection{LoRA Hyperparameters}

\begin{table}[H]
  \centering
  \caption{LoRA configuration (selected via rank ablation, see
    \cref{sec:lora:ablation}).}
  \label{tab:lora_config}
  \begin{tabular}{lcp{7cm}}
    \toprule
    \textbf{Parameter} & \textbf{Value} & \textbf{Rationale} \\
    \midrule
    Rank ($r$) & 8 & Optimal from ablation over $r \in \{2, 4, 8, 16, 32\}$ \\
    Scale ($\alpha$) & 16 & Effective scale $\alpha/r = 2.0$ \\
    Dropout & 0.1 & Regularisation \\
    Target stages & 3, 4 & High-level features; stages 0--2 frozen \\
    Target modules & QKV & All attention projections \\
    \bottomrule
  \end{tabular}
\end{table}

% ---------------------------------------------------------------------------
% 5.2 Auxiliary Semantic Heads
% ---------------------------------------------------------------------------
\subsection{Auxiliary Semantic Heads}
\label{sec:lora:aux}

To enrich the encoder's feature space for downstream SDP training, we
attach auxiliary semantic regression heads that predict tumour properties
directly from GAP-pooled encoder features:
\begin{equation}
  \hat{\mathbf{y}}_p = \pi_p\bigl(\operatorname{GAP}(\texttt{encoder10}(\mathbf{x}))\bigr),
  \quad p \in \{\text{vol}\},
  \label{eq:aux_heads}
\end{equation}
where $\pi_{\text{vol}}: \R^{768} \to \R^{1}$ is a linear head predicting
log whole-tumour volume. The auxiliary loss is:
\begin{equation}
  \mathcal{L}_{\text{aux}} = \lambda_{\text{aux}} \cdot
    \|\hat{\mathbf{y}}_{\text{vol}} - \mathbf{y}_{\text{vol}}\|_2^2.
  \label{eq:aux_loss}
\end{equation}
Following the R1 revision (\cref{sec:eval:revision}), auxiliary heads for
location and shape have been removed; only the volume head is retained. The
auxiliary loss is activated after epoch~5 with a linear warmup over 10~epochs
($\lambda_{\text{aux}} = 0.1$ at full strength).

\begin{remark}
  The auxiliary heads are discarded after Phase~1. Their sole purpose is to
  provide multi-task learning signal during encoder adaptation, encouraging
  the encoder to produce features from which tumour volume can be linearly
  decoded.
\end{remark}

% ---------------------------------------------------------------------------
% 5.3 Training Protocol
% ---------------------------------------------------------------------------
\subsection{Training Protocol}
\label{sec:lora:training}

\begin{table}[H]
  \centering
  \caption{Phase 1 training configuration.}
  \label{tab:lora_training}
  \begin{tabular}{lcp{6cm}}
    \toprule
    \textbf{Parameter} & \textbf{Value} & \textbf{Rationale} \\
    \midrule
    Epochs & 150 & With early stopping (patience 25) \\
    LR (LoRA) & $10^{-4}$ & Standard for LoRA adaptation \\
    LR (decoder) & $5 \times 10^{-4}$ & Larger for pretrained decoder \\
    LR (aux heads) & $10^{-3}$ & Small heads, fast convergence \\
    Optimizer & AdamW & Decoupled weight decay \\
    Weight decay & $10^{-5}$ & Light regularisation \\
    Scheduler & Cosine decay, 10-epoch warmup & Smooth convergence \\
    Batch size & 2 (per GPU) & Memory-constrained by $128^3$ input \\
    Gradient accumulation & 4 steps & Effective batch size 8 \\
    Gradient clipping & $\|\mathbf{g}\|_2 \leq 1.0$ & Stability \\
    Precision & bf16-mixed & Memory efficiency on A100 \\
    \bottomrule
  \end{tabular}
\end{table}

The segmentation loss combines Dice loss and cross-entropy with equal weight:
\begin{equation}
  \mathcal{L}_{\text{seg}} = \lambda_{\text{dice}}\,\mathcal{L}_{\text{Dice}}
    + \lambda_{\text{CE}}\,\mathcal{L}_{\text{CE}},
  \quad \lambda_{\text{dice}} = \lambda_{\text{CE}} = 1.0.
  \label{eq:seg_loss}
\end{equation}
The total training objective is:
\begin{equation}
  \mathcal{L} = \mathcal{L}_{\text{seg}} + w(e)\,\mathcal{L}_{\text{aux}},
  \label{eq:lora_total_loss}
\end{equation}
where $w(e) = \lambda_{\text{aux}} \cdot \min(1, \max(0, (e - 5)/10))$ is
the warmup schedule for the auxiliary loss starting at epoch~5.

% ---------------------------------------------------------------------------
% 5.4 Post-Training Protocol
% ---------------------------------------------------------------------------
\subsection{Post-Training: LoRA Merge and Checkpoint Export}
\label{sec:lora:merge}

After training, LoRA weights are absorbed into the base model weights:
\begin{equation}
  \mathbf{W}_{\text{merged}} = \mathbf{W}_0 + \frac{\alpha}{r}\,
    \mathbf{B}\mathbf{A}.
  \label{eq:lora_merge}
\end{equation}
The segmentation decoder and auxiliary heads are then discarded, yielding a
merged encoder checkpoint (\texttt{phase1\_encoder\_merged.pt}) containing
only the SwinViT and encoder1--10 weights. All encoder parameters are frozen
for subsequent phases.

% ---------------------------------------------------------------------------
% 5.5 Ablation Study
% ---------------------------------------------------------------------------
\subsection{Ablation Results}
\label{sec:lora:ablation}

\subsubsection{Rank Sweep (A1)}

We evaluate LoRA ranks $r \in \{2, 4, 8, 16, 32\}$ on segmentation Dice
and linear probe $R^2$ for volume prediction.

\begin{table}[H]
  \centering
  \caption{LoRA rank ablation (A1). Dice scores on \texttt{lora\_val} and
    linear probe $R^2$ for volume on test features.}
  \label{tab:rank_ablation}
  \begin{tabular}{cccccc}
    \toprule
    \textbf{Rank} & \textbf{LoRA params} & \textbf{Dice (WT)} &
      \textbf{Dice (TC)} & \textbf{Dice (ET)} & \textbf{Vol $R^2$} \\
    \midrule
    Frozen & 0 & \textit{pending} & \textit{pending} & \textit{pending} & \textit{pending} \\
    2  & $\sim$49K  & \textit{pending} & \textit{pending} & \textit{pending} & \textit{pending} \\
    4  & $\sim$99K  & \textit{pending} & \textit{pending} & \textit{pending} & \textit{pending} \\
    8  & $\sim$197K & \textit{pending} & \textit{pending} & \textit{pending} & \textit{pending} \\
    16 & $\sim$394K & \textit{pending} & \textit{pending} & \textit{pending} & \textit{pending} \\
    32 & $\sim$787K & \textit{pending} & \textit{pending} & \textit{pending} & \textit{pending} \\
    \bottomrule
  \end{tabular}
\end{table}

\begin{figure}[H]
  \centering
  \fbox{\parbox{0.8\textwidth}{\centering\vspace{3cm}
    \textbf{Figure placeholder:} LoRA rank ablation curves.\\
    Left: Dice vs.\ rank. Right: Probe $R^2$ vs.\ rank.
  \vspace{3cm}}}
  \caption{LoRA rank ablation (A1). Left panel shows validation Dice scores
    per sub-region as a function of LoRA rank. Right panel shows linear probe
    $R^2$ for volume prediction. Rank $r=8$ achieves the best trade-off
    between capacity and parameter efficiency.}
  \label{fig:rank_ablation}
\end{figure}

\subsubsection{LoRA vs.\ DoRA (A2)}

DoRA (\cref{eq:dora}) is evaluated as an alternative adapter type at $r=8$.

\begin{table}[H]
  \centering
  \caption{LoRA vs.\ DoRA comparison (A2) at rank 8.}
  \label{tab:lora_vs_dora}
  \begin{tabular}{lccc}
    \toprule
    \textbf{Adapter} & \textbf{Dice (WT)} & \textbf{Vol $R^2$ (linear)} & \textbf{Vol $R^2$ (RBF)} \\
    \midrule
    LoRA ($r=8$) & \textit{pending} & \textit{pending} & \textit{pending} \\
    DoRA ($r=8$) & \textit{pending} & \textit{pending} & \textit{pending} \\
    \bottomrule
  \end{tabular}
\end{table}

% ---------------------------------------------------------------------------
% 5.6 Dual-Domain Adaptation
% ---------------------------------------------------------------------------
\subsection{Dual-Domain Adaptation}
\label{sec:lora:dual}

A key scientific question is whether a single LoRA-adapted encoder can produce
semantically structured features for both glioma and meningioma simultaneously.
Dual-domain training combines BraTS-MEN ($\sim$750~studies) and BraTS-GLI
($\sim$950~scans) via a \texttt{ConcatDataset} with domain-balanced
\texttt{WeightedRandomSampler} (50/50 domain weight).

\subsubsection{Empirical Findings}

\begin{table}[H]
  \centering
  \caption{Single-domain vs.\ dual-domain LoRA adaptation results.}
  \label{tab:dual_domain}
  \begin{tabular}{lcccccc}
    \toprule
    & \textbf{MEN Dice} & \textbf{Vol $R^2$} & \textbf{CV Vol $R^2$}
    & \textbf{MMD$^2$} & \textbf{CKA} & \textbf{erank} \\
    \midrule
    Frozen baseline & 0.52 & 0.56 & --- & 0.118 & 0.002 & 47.6 \\
    MEN LoRA $r$=8  & 0.48 & 0.57 & --- & 0.116 & 0.002 & 45.8 \\
    Dual LoRA $r$=8 & 0.57 & 0.62 & 0.712 & 0.037 & 0.017 & 25.4 \\
    \bottomrule
  \end{tabular}
\end{table}

The dual-domain encoder achieves the best volume probe $R^2$ and reduces the
domain gap (MMD$^2$: $0.118 \to 0.037$, a 68.6\% reduction), but at the cost
of a substantial drop in effective rank ($47.6 \to 25.4$, a 46.6\% reduction),
indicating representation compression.

\subsubsection{Critical Finding: Cross-Domain Transfer Failure}

Despite the reduced domain gap metrics, cross-domain GP probes reveal a
critical limitation: training probes on GLI features and evaluating on MEN
features yields $R^2 \approx -0.01$ for volume, indicating that the shared
feature space does not support cross-domain semantic transfer. Furthermore,
MEN segmentation Dice degrades from 0.64 to 0.36 under dual training,
indicating catastrophic interference rather than beneficial transfer.

\begin{remark}
  The dual-domain results provide important negative evidence: domain gap
  reduction (measured by MMD, CKA) does not guarantee cross-domain semantic
  compatibility. The effective rank drop from 47.6 to 25.2 suggests that the
  encoder resolves the domain conflict by compressing both domains into a
  lower-dimensional shared subspace that loses domain-specific semantic
  content.
\end{remark}

\subsubsection{LoRA Marginality}

A surprising finding is the marginal benefit of LoRA adaptation for semantic
probes: the frozen baseline achieves volume $R^2 = 0.794$, while the best
LoRA condition improves this by only $\Delta R^2 = +0.016$ (within noise).
This suggests that BrainSegFounder's SSL pretraining already produces
features with reasonable meningioma volume decodability, and the primary
benefit of LoRA is in segmentation quality rather than downstream feature
quality.

% ---------------------------------------------------------------------------
% 5.7 Configuration
% ---------------------------------------------------------------------------
\subsection{Configuration}
\label{sec:lora:config}

All LoRA experiments are configured via YAML files under
\texttt{experiments/lora/config/}. Server configurations use 250/50/200/500
splits (lora\_train/lora\_val/probe\_train/test), differing slightly from the
canonical specification. Four conditions are defined:
LoRA/DoRA $\times$ semantic/no-semantic heads.
